
TLCM connects to several open questions in feature-based interpretability. Under the propose-and-reject hypothesis, the residual stream at any layer $i$ contains both persistent features (those that have survived correction from prior layers) and uncorrected proposals from layer $i$ that layer $i+1$ may reverse. Methods that decompose the residual stream at a single layer cannot distinguish between these two categories. We highlight three connections below.

\textbf{Feature descriptions lack high specificity.} SAEs trained on a single layer's output cannot separate persistent features from soon-to-be-reversed proposals, predicting \textit{misfires}: features that activate on content the model is about to correct. Consistent with this, recent work \cite{templeton2024scaling} observed that over 50\% of activated features from an SAE are labeled by a large LLM as ``Irrelevant'' or ``Only vaguely related'' to the text on which they fire.

\textbf{Effective model steering may need to overcome TLCM's correction.} If TLCM corrections oppose amplified features, then steering interventions may need to exceed the correction capacity to take effect. As shown in Appendix \ref{appendix:tlcm_adaptivity}, TLCM's correction capacity begins to diminish when the prior layer is amplified beyond $2\times$. Recent work found that effective steering requires extreme amplification \textbf{}levels, up to $10\times$ a feature's maximum observed value \cite{templeton2024scaling}, which may reflect the need to overcome correction mechanisms. We propose that intervening on feature directions that TLCM does not target, identified using the Jacobian, could be a promising direction for future work.

\textbf{Transcoders outperform SAEs.} Recent work finds that per-layer transcoders achieve both more faithful reconstruction and more interpretable features than SAEs \cite{paulo2025transcodersbeatsparseautoencoders}. TLCM offers one interpretation of this gap. Unlike SAEs, transcoders model a layer's computation as a function of its input residual stream, so they can in principle represent corrections as input-conditional mechanisms that distinguish features being corrected from features being newly contributed. The cross-layer architecture \cite{ameisen2025circuit} additionally captures dynamics across adjacent layers: a CLT feature that encodes at layer $i$ and decodes to multiple downstream layers can represent both the proposal and its subsequent correction.

Whether TLCM is the most efficient mechanism available, or a consequence of architectural constraints, is an open question. Architectural modifications that reduce TLCM during training, if they lead to lower loss, would suggest the latter. It is also unknown whether TLCM develops in SSM-based sequence models or transformer variations (gated attention, NoPE, MOE).

While the propose-and-reject hypothesis is consistent with our results, it remains a conceptual framework without formalization. One strategy towards validating it might be developing techniques to decompose the correction into semantically meaningful concepts. By employing recent circuit tracing and automated feature labeling techniques, we might be able to semantically explain the correction and its causes.

As a well-characterized and broadly observed phenomenon, TLCM may serve as a useful constraint for researchers building mechanistic accounts of transformer computation.

\newpage
